% Created 2025-02-05 wo 09:23
% Intended LaTeX compiler: pdflatex
\documentclass[11pt]{article}
\usepackage[utf8]{inputenc}
\usepackage[T1]{fontenc}
\usepackage{graphicx}
\usepackage{longtable}
\usepackage{wrapfig}
\usepackage{rotating}
\usepackage[normalem]{ulem}
\usepackage{amsmath}
\usepackage{amssymb}
\usepackage{capt-of}
\usepackage{hyperref}
\author{Music Player Daemon (MPD) user}
\date{\today}
\title{pop-test manual}
\hypersetup{
 pdfauthor={Music Player Daemon (MPD) user},
 pdftitle={pop-test manual},
 pdfkeywords={},
 pdfsubject={},
 pdfcreator={Emacs 29.4 (Org mode 9.6.15)}, 
 pdflang={English}}
\begin{document}

\maketitle
\tableofcontents

In Dutch 🇳🇱, \emph{pop} means puppet or doll. But pop-test is not about dolls! It’s a tool that helps test computer programs in an easy and future-proof manner that improves reliability and engineering freedom.

pop-test simplifies end-to-end testing by providing a powerful orchestration tool and expressive language for your applications and their dependencies.

Define your test scenarios in a DSL (YAML, in the future also TOML or JSON) and automate the spinning up (and down) of complex environments involving multiple programs, databases, and services, executing comprehensive tests with ease.

pop-test is built with Rust, enhancing safety, performance and reliability, with a modern and accessible code-base.

pop-tests as we affectionately call them, can range from small scope integration tests to all encompassing end-to-end or black box tests.


\section{Why pop-test?}
\label{sec:org50786cf}

\begin{itemize}
\item Catches mistakes early – Helps find problems before real users do.
\item Test coverage with a real meaning - test your systems from head to toe
\item Frees you and allows you to refactor and test your systems more easily
\item Saves time – No need to test every little part manually, test on higher level.
\item Easy to use – Just write simple rules, in a comfy DSL and \texttt{pop-test} does the rest!
\item Fast and reliable – Built with Rust, a programming language known for speed and safety.
\end{itemize}

\section{Motivation}
\label{sec:org3879632}

Imagine you're building a toy robot, but you don't need to understand all the small parts inside the robot to know if it works. You just need to press a button and see if it performs the task correctly.

A big advantage of this is that you test based on the output, not the code itself or implementation.

Of course this doesn't mean you shouldn't add some unit/property-based tests when useful. But these shouldn't be the bulk of your tests.

It's great when you change or improve your code (refactor), that you don’t need to worry about breaking anything as long as the system gives the same results as before.

It's kind of like saying: "As long as pressing the button still works the same, everything’s fine."

These tests are a safety net for changes to systems (specially if you aren't too familiar with the system yet). It gives you confidence when making changes to the code.

How Does This "Free You"? You can change the internal code without fear because you already have pop-tests that check if everything works as expected. So, you can keep improving the system without worrying about accidentally breaking things.

This also allows for faster development, since you don’t have to check every single little part of the code manually.

As your system grows, you don’t need to rewrite tests every time. You just keep running your pop-tests, and you know that new features won’t break old ones.

In short, pop-tests make engineers happier, business people more certain of what they sell and overall, it makes refactoring and maintaining things easier.

\section{The DSL (Domain Specific Language)}
\label{sec:org7cf57b9}

In order to provide an expressive and powerful way to describe test scenarios, a lot of care has gone into creating a good domain-specific language, which remains flexible but concise.

Here follows an explanation of every individual field and parameter possible, including default or optional ones, as well as some recommendations.

\section{Example tests, using the DSL}
\label{sec:orga493b21}

Here follow some examples of pop-tests.

In this example we spin up a REST API written in Scala, with SBT, and poll its \texttt{/health} endpoint for 90 seconds or until healthy. Then we perform some checks on the response of certain HTTP endpoints. These checks are performed using \texttt{jql} and this allows us to check deeply nested attributes of the JSON response.

\begin{verbatim}
pop_test:
  pop_services:
    - id: "green-energy-mix"
      dir: "$HOME/Ontwikkeling/Werk"
      base_url: "http://localhost:50000"
      start_cmd_args:
        - 'sbt "greenEnergyMix/run"'
      env_vars:
        SOME_TOKEN: "0027809ce9"
      health_check:
        path: "/health"
        patience_seconds: 90

  pop_checks:
    - check_url: "/api/v1/forecast"
      service_id: "green-energy-mix"
      json_body_patterns:
        - pattern: '"meta""forecastDate"'
          condition: '!= ""'
        - pattern: '"data"[0]"load"'
          condition: "> 0"
        - pattern: '"data"[0]"solar"'
          condition: ">= 0"

\end{verbatim}
\end{document}
